\chapter{Paradoxos de Representação e Formalização}
\label{chap:apendiceRepresentacaoFormalizacao}

Adicionar Introdução

\section{O Paradoxo de Skolem (1922)}

Apresentar o Teorema de Löwenheim-Skolem.

\section{O Paradoxo de König (1905)}

Apresentar o episódio histórico no Congresso Internacional de
Matemática de 1904.

\section{O Teorema de Goodstein (1944)}

Apresentar a sequência de Goodstein e seu crescimento explosivo
seguido de colapso inevitável a zero. Demonstrar a prova via
ordinais transfinitos: associar a cada termo um ordinal em
notação hereditária com \(\omega\) no lugar da base e usar a
boa-ordenação dos ordinais para concluir. Conectar com o
Capítulo~7 (Números Ordinais). Apresentar o resultado de
Kirby e Paris (1982): a terminação não é demonstrável dentro
da aritmética de Peano,
situando o teorema na mesma família do resultado de Gödel
do Apêndice~B, mas com enunciado elementar e papel dos
ordinais completamente explícito.